\section{Métodos topológicos de existencia: disparo, puntos fijos y grado}
\label{sec:metodos-topologicos-existencia}

\HAFOCardExistence

Los métodos topológicos garantizan la existencia de ceros o puntos fijos a
partir de propiedades de continuidad, compacidad y comportamiento en la
frontera de un dominio \cite{Amster2021MetodosTopologicos}. Aplicados a un
campo o a un mapa de retorno, ayudan a localizar equilibrios y ciclos y a
demostrar su persistencia. La estabilidad y la ocultedad requieren, además,
un análisis dinámico y de cuencas.

\subsection{De una ecuación diferencial a un operador}

Hay tres traducciones fundamentales:

\begin{enumerate}
 \item un equilibrio es un cero de $f$:
 \[
 f(e)=0;
 \]
 \item al elegir una sección local transversal en un punto de una órbita
 periódica, el retorno local asociado tiene allí un punto fijo:
 \[
  P(z)=z;
 \]
 en una sección global con varias intersecciones, la misma órbita puede
 corresponder a un punto periódico de periodo \(k>1\);
 \item una solución del PVI autónomo de Caputo
 \[
  {}^{\mathrm C}D_{t_0+}^{q}x(t)=f(x(t)),\qquad
  x(t_0)=x_0,\qquad 0<q<1,
 \]
 es un punto fijo del operador de Volterra:
 \begin{equation}
  (\mathcal Kx)(t)=x_0+
  \frac{1}{\Gamma(q)}\int_{t_0}^{t}(t-\tau)^{q-1}
  f(x(\tau))\,d\tau.
  \label{eq:guide-caputo-fixed-operator}
 \end{equation}
\end{enumerate}

El primer problema es finito-dimensional; el tercero vive en un espacio de
funciones. El segundo puede ser finito-dimensional, pero exige que el primer
retorno esté definido y dependa continuamente de la condición inicial.

La interpretación de la tercera equivalencia requiere fijar el dominio
del operador fraccionario. Para una función continua $x$ con $x(t_0)=x_0$,
la derivada regularizada se define por
\begin{equation}
 {}^{\mathrm C}D_{t_0+}^{q}x
 :=\frac{d}{dt}I_{t_0+}^{1-q}(x-x_0),
 \label{eq:guide-regularized-caputo}
\end{equation}
cuando la integral del miembro derecho es absolutamente continua.
Si $x=x_0+I^qg$ con $g=f\circ x$ continua, la composición de integrales
da $I^{1-q}(x-x_0)=I^1g$ y, al derivar, se obtiene $g$.
La composición se justifica por Fubini: $g$ es acotada en el intervalo,
y la integral interior del producto de núcleos es
\[
 \int_s^t(t-u)^{-q}(u-s)^{q-1}\,du
 =B(1-q,q)=\Gamma(1-q)\Gamma(q).
\]
Sus singularidades son integrables porque $0<q<1$. Recíprocamente, si
\cref{eq:guide-regularized-caputo} vale $g$, se integra para obtener
$I^{1-q}(x-x_0)=I^1g$, ya que ambas expresiones se anulan en $t_0$.
Aplicar $I^q$ da $I^1(x-x_0)=I^{q+1}g$; derivar recupera
$x-x_0=I^qg$. Si además $x$ es absolutamente continua, escribir
$x-x_0=I^1x'$ muestra que
\cref{eq:guide-regularized-caputo} coincide con la definición clásica
$I^{1-q}x'$. La existencia obtenida abajo con $f$ meramente continua
es, por tanto, existencia de una solución de Volterra y de esta derivada
regularizada; usar la definición clásica exige verificar adicionalmente
la continuidad absoluta de $x$.

\subsection{Espacio funcional mínimo para los argumentos de existencia}

En el intervalo compacto \(J=[t_0,T]\), el espacio
\begin{equation}
 X=C(J,\R^n),\qquad
 \|x\|_\infty=\max_{t\in J}\|x(t)\|,
 \label{eq:guide-sup-norm}
\end{equation}
es completo: toda sucesión de Cauchy uniforme converge a una función continua.
Por tanto, \(X\) es un espacio de Banach. Esta propiedad es la que permite
aplicar el teorema de contracción al operador integral, no sólo el hecho de que
sus elementos sean funciones.

Un operador \(K:X\to X\) es \emph{compacto y continuo} cuando
es continuo y envía conjuntos acotados en conjuntos relativamente compactos.
Para operadores no lineales, la segunda propiedad no implica por sí sola la
primera. En \(C(J,\R^n)\), el criterio de
Arzelà--Ascoli reduce esta comprobación a dos tareas: acotación uniforme y
equicontinuidad de la familia imagen. El núcleo de Volterra proporciona un
módulo de continuidad de orden fraccionario, de modo que puede hacer compacto
al operador de Caputo aun cuando éste no sea contractivo. Compacidad y
contracción cumplen funciones distintas: la primera apoya existencia por grado;
la segunda añade unicidad e iteración convergente
\cite{Deimling1985,Amster2021MetodosTopologicos}.

\paragraph{Ejercicio sobre compacidad.}
Demuestre que una bola cerrada de \(C(J,\R^n)\) no es compacta en general;
después identifique qué hipótesis adicional de Arzelà--Ascoli falta. Éste es el
contraejemplo que impide sustituir ``acotado'' por ``compacto'' en dimensión
infinita.

\subsection{Disparo: transformar condiciones de frontera en un cero}

Considérese
\begin{equation}
 u''=F(t,u,u'),\qquad u(0)=a,\qquad u(T)=b.
 \label{eq:guide-shooting-bvp}
\end{equation}
Para cada pendiente inicial $s$, resuélvase el PVI
$u(0)=a$, $u'(0)=s$ y defínase
\begin{equation}
 S(s)=u(T;s)-b.
 \label{eq:guide-shooting-residual}
\end{equation}
Entonces $S(s_*)=0$ si y sólo si el disparo $s_*$ resuelve el problema de
frontera. La equivalencia es inmediata; su aplicación requiere demostrar que
$S$ está definido y es continuo en un intervalo adecuado.
Una condición suficiente es que $F$ sea continua en $t$, localmente
Lipschitz en $(u,u')$ y que las soluciones correspondientes a las
pendientes del intervalo permanezcan en un compacto común hasta $T$.
La teoría del PVI da entonces unicidad, existencia hasta $T$ y dependencia
continua respecto de $s$; evaluar en $T$ conserva esa continuidad.

El análogo para un ciclo es un residual de retorno. Si $z\in\Sigma$ y
$\tau(z)>0$ es el primer tiempo de retorno,
\begin{equation}
 P(z)=\varphi_{\tau(z)}(z),\qquad R(z)=P(z)-z.
 \label{eq:guide-return-residual}
\end{equation}
Un cero de $R$ proporciona un ciclo. En el PLL, el retorno se construye sobre
el cilindro; en MAVPD, una sección transversal reduce localmente el problema
de una órbita en $\R^3$ a la iteración de un mapa en dimensión dos.

Para justificar la continuidad local del retorno, escriba la sección como
$\Sigma=\{h=0\}$, con $h$ de clase $C^1$ y $Dh\ne0$.
Si $\varphi_{\tau_*}(z_*)\in\Sigma$ y
$Dh(\varphi_{\tau_*}(z_*))f(\varphi_{\tau_*}(z_*))\ne0$, el teorema
de la función implícita aplicado a
$h(\varphi_t(z))=0$ produce una rama $C^1$ de tiempos $\tau(z)$,
siempre que $f\in C^1$. Para identificarla como primer retorno debe
excluirse un cruce anterior, además de asegurar existencia hasta ese tiempo.
Sobre una región donde estas condiciones persisten, $P$ es $C^1$.
Al derivar la identidad de cruce, con $n^{\mathsf T}=Dh(P(z))$, se obtiene
\begin{align}
 D\tau(z)&=-\frac{n^{\mathsf T}D_z\varphi_{\tau(z)}(z)}
                   {n^{\mathsf T}f(P(z))},\nonumber\\
 DP(z)&=\left[I-\frac{f(P(z))n^{\mathsf T}}
                        {n^{\mathsf T}f(P(z))}\right]
                        D_z\varphi_{\tau(z)}(z).
 \label{eq:guide-return-derivative}
\end{align}
En estas fórmulas, $D_z\varphi_{\tau(z)}$ deriva el flujo manteniendo
fijo el tiempo de evaluación; la contribución de $D\tau$ aparece en
el factor de proyección. La restricción a $T_z\Sigma$ es la derivada
del retorno en la sección. Así se separa la estabilidad transversal
del movimiento a lo largo de la órbita.

\subsection{Teorema de contracción: existencia, unicidad y aproximación}

\begin{proposition}[punto fijo de Banach]
Sea $(X,d)$ un espacio completo no vacío y $T:X\to X$ una contracción: existe $0<L<1$ tal que
$d(Tx,Ty)\leq Ld(x,y)$. Entonces $T$ posee un único punto fijo $x_*$ y, para
cualquier $x_0$, la iteración $x_{n+1}=Tx_n$ converge a $x_*$.
\end{proposition}

\begin{proof}
La estimación
\[
 d(x_{n+1},x_n)\leq L^n d(x_1,x_0)
\]
y, para \(m>n\), la desigualdad triangular dan
\[
 d(x_m,x_n)\le\sum_{j=n}^{m-1}L^j d(x_1,x_0)
 \le\frac{L^n}{1-L}d(x_1,x_0).
\]
El último miembro tiende a cero, por lo que $(x_n)$ es de Cauchy. Por
completitud, $x_n\to x_*$. La continuidad de $T$ da $Tx_*=x_*$. Si $y_*$ fuera
otro punto fijo,
$d(x_*,y_*)\leq Ld(x_*,y_*)$, de donde $x_*=y_*$.
\end{proof}

Para aplicar el teorema, elija una bola de estados
\(B=\overline B(x_0,R)\) donde \(f\) sea \(L\)-Lipschitz, y sea
\(M=\max_{u\in B}\|f(u)\|\). En el espacio de funciones use la bola
cerrada completa
\[
 \mathcal D=\{x\in C([t_0,T],\R^n):\|x-x_0\|_\infty\le R\},
 \qquad \Delta=T-t_0,
\]
donde \(x_0\) también designa la función constante. Para \(x,y\in\mathcal D\),
integrar el núcleo da
\begin{align*}
 \|\mathcal Kx-x_0\|_\infty
 &\le\frac{M\Delta^q}{\Gamma(q+1)},\\
 \|\mathcal Kx-\mathcal Ky\|_\infty
 &\le\frac{L\Delta^q}{\Gamma(q+1)}\|x-y\|_\infty,
\end{align*}
pues \(\int_{t_0}^t(t-s)^{q-1}\,ds=(t-t_0)^q/q\) y
\(q\Gamma(q)=\Gamma(q+1)\). Si
\(M\Delta^q/\Gamma(q+1)\le R\) y
\(L\Delta^q/\Gamma(q+1)<1\), entonces
\(\mathcal K:\mathcal D\to\mathcal D\) es una contracción.
Ambas condiciones se cumplen al tomar \(\Delta\) suficientemente pequeño.
Esto prueba existencia y unicidad local del PVI de Caputo: cualquier solución
con el mismo dato permanece inicialmente en \(B\), por continuidad, y
coincide allí con el punto fijo. No prueba existencia de un atractor ni
periodicidad. Para
 \(m-1<q<m\) se requieren \(m\) datos iniciales y el término libre incluye el
 polinomio inicial correspondiente.

\subsection{Brouwer y el grado en dimensión finita}

\begin{proposition}[punto fijo de Brouwer]
Toda aplicación continua $T:\overline B\to\overline B$ de una bola cerrada de
$\R^n$ en sí misma tiene al menos un punto fijo.
\end{proposition}

La conclusión es existencial: no afirma unicidad, estabilidad ni proporciona
un algoritmo rápido. El grado de Brouwer refina esta idea cuando
$\Omega\subset\R^n$ es abierta y acotada, el mapa
$F:\overline\Omega\to\R^n$ es continuo y
$y\notin F(\partial\Omega)$. Se denota
\[
 \deg(F,\Omega,y)\in\mathbb Z.
\]
Sus propiedades operativas son:

\begin{enumerate}
 \item \textbf{Existencia:} si $\deg(F,\Omega,y)\neq0$, entonces
 $F(x)=y$ para algún $x\in\Omega$.
 \item \textbf{Normalización:}
 $\deg(I,\Omega,y)=1$ cuando $y\in\Omega$.
 \item \textbf{Aditividad:} el grado se suma al separar regiones que contienen
 todos los ceros.
 \item \textbf{Invariancia homotópica:} si
 $H:[0,1]\times\overline\Omega\to\R^n$ es continua y
 $y\notin H(\lambda,\partial\Omega)$ para todo $\lambda$, entonces
 \begin{equation}
  \deg(H(0,\cdot),\Omega,y)
  =\deg(H(1,\cdot),\Omega,y).
  \label{eq:guide-degree-homotopy}
 \end{equation}
\end{enumerate}

Si $F$ es $C^1$, $y$ es valor regular y hay un número finito de preimágenes,
\begin{equation}
 \deg(F,\Omega,y)=
 \sum_{x\in F^{-1}(y)\cap\Omega}
 \operatorname{sgn}\det DF(x).
 \label{eq:guide-degree-regular}
\end{equation}
Esta fórmula expresa el grado como la suma de las orientaciones locales de
los ceros regulares.

\paragraph{Ejemplo.}
Para $F(x)=x-x^3$ en $\Omega=(-2,2)$, los ceros son $-1,0,1$ y
$F'(x)=1-3x^2$. Por tanto,
\[
 \deg(F,(-2,2),0)=-1+1-1=-1.
\]
El grado no nulo garantiza que al menos un equilibrio persiste bajo toda
homotopía que no haga cruzar un cero por la frontera. El cálculo no determina
cuáles de los tres equilibrios son atractores; esa información requiere el
signo de la linealización o un argumento dinámico.

\subsection{Grado del residual de retorno y continuación}

Para ciclos se aplica el grado a
\begin{equation}
 F(z)=P(z)-z.
 \label{eq:guide-degree-return}
\end{equation}
Sea \(\Omega\) un abierto acotado en una carta de la sección, y supóngase
que \(P\) está definido y es continuo en \(\overline\Omega\), con
\(P(z)\ne z\) en \(\partial\Omega\). Si
$\deg(P-I,\Omega,0)\neq0$, existe un punto fijo de $P$ en $\Omega$ y, por
consiguiente, una órbita periódica.
En una familia $P_\lambda$, la invariancia homotópica conserva existencia
mientras ningún cero alcance $\partial\Omega$ y el retorno no pierda su
dominio. Esto proporciona una base topológica para la continuación, distinta
de simplemente reutilizar numéricamente el último estado.

La continuación armónica proporciona una semilla aproximada. La resolución
de $P_\lambda(z)-z=0$ corrige esa aproximación, y el cálculo validado del
grado o del índice local establece existencia en un dominio. Las tres etapas
son:

\[
 \text{predictor armónico}
 \longrightarrow
 \text{corrector de retorno}
 \longrightarrow
 \text{existencia demostrada en la región}.
\]

\subsection{Ejemplo resuelto: localización y clasificación de un ciclo oculto}
\label{subsec:guide-hidden-cycle-topology}

Considérese el sistema polinomial planar
\begin{equation}
 \begin{aligned}
 \dot x&=-Q(x,y)x-y,\qquad
 \dot y=x-Q(x,y)y,\\
 Q(x,y)&=(x^2+y^2-a^2)(x^2+y^2-b^2),
 \end{aligned}
 \label{eq:guide-hidden-cycle-field}
\end{equation}
con \(0<a<b\). Para \(r=\sqrt{x^2+y^2}>0\), se calcula
\[
 \dot r=\frac{x\dot x+y\dot y}{r}
 =\frac{-Q(x^2+y^2)}r=-Qr,\qquad
 \dot\theta=\frac{x\dot y-y\dot x}{r^2}
 =\frac{x^2+y^2}{r^2}=1.
\]
Por tanto,
\begin{equation}
 \dot r=g(r):=-r(r^2-a^2)(r^2-b^2),
 \qquad \dot\theta=1.
 \label{eq:guide-hidden-cycle-radial}
\end{equation}
El único equilibrio es el origen, pues la matriz
\(\left(\begin{smallmatrix}-Q&-1\\1&-Q\end{smallmatrix}\right)\)
que multiplica a \((x,y)\) tiene determinante \(Q^2+1>0\).
En el origen, \(Q=a^2b^2\), y los términos procedentes de derivar
\(Q\) están multiplicados por \(x\) o \(y\) y se anulan.
El polinomio característico del Jacobiano es
\((\lambda+a^2b^2)^2+1\). Sus autovalores son
\(-a^2b^2\pm\mathrm i\), de modo que es un foco asintóticamente estable.
Los círculos \(C_a=\{r=a\}\) y \(C_b=\{r=b\}\) son órbitas periódicas
de periodo \(2\pi\).

El signo de \(g\) determina todos los destinos radiales:
\[
 \begin{array}{c|ccc}
 r & (0,a)&(a,b)&(b,\infty)\\ \hline
 \dot r&<0&>0&<0
 \end{array}
\]
Las soluciones están acotadas hacia adelante y
\begin{equation}
 \mathcal B(\{0\})=\{r<a\},\qquad
 \mathcal B(C_b)=\{r>a\},\qquad
 \partial\mathcal B(C_b)=C_a.
 \label{eq:guide-hidden-cycle-basins}
\end{equation}
En efecto, la unicidad impide cruzar los radios invariantes \(a\) y \(b\),
y \(r(t)\le\max\{r(0),b\}\). Esta cota y la regularidad polinómica
excluyen explosión en tiempo finito. En cada intervalo de la tabla,
\(r(t)\) es monótona y acotada, por lo que tiene un límite \(\ell\).
Si \(g(\ell)\ne0\), por continuidad \(\dot r\) conservaría un signo y
un módulo separado de cero para tiempos grandes, lo cual impediría ese
límite. Así, \(g(\ell)=0\). La dirección de monotonía selecciona
\(r(t)\to0\) si \(r(0)<a\), y \(r(t)\to b\) si \(r(0)>a\).
Como la distancia de un punto al círculo \(C_b\) es \(|r-b|\), esto
prueba las cuencas indicadas. Su frontera común es \(C_a\), una
separatriz repulsiva, y \(C_b\)
es un ciclo atractivo. Toda bola centrada en el origen de radio menor que
\(a\) queda fuera de \(\mathcal B(C_b)\); el ciclo \(C_b\) es oculto.

La misma geometría permite localizar el ciclo por retorno. Elija
\(a<r_-<b<r_+\). En el anillo
\(N=\{r_-\leq r\leq r_+\}\), el campo apunta hacia dentro en ambas
componentes de la frontera. Sobre la semirrecta \(\theta=0\), el retorno
\(P\) tiene tiempo constante \(2\pi\) y satisface
\[
 P(r_-)-r_->0,\qquad P(r_+)-r_+<0.
\]
Para calcular el grado, defina
\(H_\lambda(r)=(1-\lambda)[P(r)-r]+\lambda(b-r)\).
Ambos sumandos son positivos en \(r_-\) y negativos en \(r_+\);
por tanto, \(H_\lambda\) nunca se anula en la frontera del intervalo.
La invariancia homotópica y la derivada \((b-r)'=-1\) dan
\[
 \deg(P-I,(r_-,r_+),0)
 =\deg(b-I,(r_-,r_+),0)=-1.
\]
Existe así un punto fijo. El análisis radial lo identifica como \(r=b\):
fuera de ese radio, \(r(t)\) es estrictamente monótona y no puede volver
a su valor inicial en tiempo \(2\pi\).
La derivada \(u(t)=\partial r(t;r_0)/\partial r_0\) satisface
\(\dot u=g'(r(t;r_0))u\), \(u(0)=1\). En el ciclo \(r=b\),
\(u(t)=e^{g'(b)t}\). Además,
\[
 g'(r)=-5r^4+3(a^2+b^2)r^2-a^2b^2,\qquad
 g'(b)=-2b^2(b^2-a^2).
\]
Su multiplicador transversal es, por ello,
\begin{equation}
 P'(b)=\exp\bigl(2\pi g'(b)\bigr)
 =\exp\bigl[-4\pi b^2(b^2-a^2)\bigr]\in(0,1).
 \label{eq:guide-hidden-cycle-multiplier}
\end{equation}
Análogamente, \(g'(a)=2a^2(b^2-a^2)>0\), de modo que
\(P'(a)=\exp[4\pi a^2(b^2-a^2)]>1\), lo que confirma la
repulsión transversal de \(C_a\).
El grado localiza un ciclo, el multiplicador establece su estabilidad y la
frontera \(C_a\) demuestra la separación de su cuenca respecto del
equilibrio. Son argumentos complementarios: el grado del campo en cualquier
disco centrado en el origen es \(+1\), porque el único cero es regular y
\(\det Df(0)=a^4b^4+1>0\). Sólo detecta ese equilibrio; no
detecta los dos ciclos. El atractor oculto del ejemplo es periódico, lo que
muestra que ocultedad y caos son propiedades distintas.

\begin{hafocommonfigure}[H]
\centering
\resizebox{0.94\linewidth}{!}{%
\begin{tikzpicture}[font=\small,>=Latex]
  \begin{scope}[shift={(2.7,2.75)}]
    \fill[VerdeClaro] (-2.5,-2.5) rectangle (2.5,2.5);
    \fill[AzulClaro] (0,0) circle (1);
    \draw[->,GrisOscuro] (-2.55,0)--(2.6,0) node[right] {$x$};
    \draw[->,GrisOscuro] (0,-2.55)--(0,2.6) node[above] {$y$};
    \draw[Rojo,very thick,dashed] (0,0) circle (1);
    \draw[Verde,very thick] (0,0) circle (2);
    \draw[->,Verde,very thick] (20:2) arc (20:65:2);
    \draw[->,Verde,very thick] (200:2) arc (200:245:2);
    \draw[->,Rojo,thick] (115:1) arc (115:155:1);
    \fill[Azul] (0,0) circle (0.055);
    \node[Azul,below left] at (0,0) {$0$};
    \node[fill=AzulClaro,text=Azul] at (-0.07,-0.58) {$r<1$};
    \node[Rojo,fill=white,inner sep=1.4pt] at (0.75,0.9) {$C_1$};
    \node[Verde,fill=VerdeClaro,inner sep=1.4pt] at (-1.66,1.35) {$C_2$};
    \node[Verde] at (1.65,-2.23) {$r>1$};
  \end{scope}
  \node[font=\bfseries,anchor=west] at (5.95,5.25)
    {Dinámica radial};
  \node[anchor=west] at (5.95,4.68)
    {$\dot r=-r(r^2-1)(r^2-4)$};
  \draw[->,GrisOscuro] (6.05,3.8)--(12.75,3.8) node[right] {$r$};
  \foreach \x/\lab in {6.2/0,8.35/1,10.5/2}{
    \draw (\x,3.68)--(\x,3.92);
    \node[below] at (\x,3.63) {$\lab$};
  }
  \draw[->,Azul,very thick] (7.85,4.12)--(6.65,4.12);
  \draw[->,Verde,very thick] (8.8,4.12)--(10.05,4.12);
  \draw[->,Verde,very thick] (12.25,4.12)--(11,4.12);
  \node[anchor=west,Azul] at (5.95,2.76)
    {$\mathcal B(\{0\})=\{r<1\}$};
  \node[anchor=west,Verde] at (5.95,2.08)
    {$\mathcal B(C_2)=\{r>1\}$};
  \node[anchor=west,Rojo] at (5.95,1.40)
    {$\partial\mathcal B(C_2)=C_1$};
  \node[anchor=west] at (5.95,0.62)
    {$B_{1/2}(0)\cap\mathcal B(C_2)=\varnothing$};
\end{tikzpicture}%
}
\caption{Cuencas exactas del sistema \eqref{eq:guide-hidden-cycle-field}
con \(a=1\) y \(b=2\). El disco azul converge al origen; la región verde
converge al ciclo \(C_2\). El ciclo repulsivo \(C_1\) separa ambas cuencas.
Las flechas radiales indican el signo de \(\dot r\), mientras las flechas
circunferenciales representan \(\dot\theta=1\).}
\label{fig:ciclo-oculto-exacto}
\end{hafocommonfigure}


\ifHAFOPedagogy
\paragraph{Ejercicio.}
Fije \(a=1\) y \(b=2\). Calcule los multiplicadores de ambos ciclos y
los destinos de reinicios con radios \(1/2\), \(3/2\) y \(3\).
Explique qué parte de la clasificación seguiría siendo desconocida si sólo
se dispusiera del índice del origen o del grado del retorno en el anillo.
\fi

\subsection{Leray--Schauder y el operador integral de Caputo}

En un espacio de Banach $X$, sea $\Omega\subset X$ abierta y acotada, y sea
$\mathcal K:\overline\Omega\to X$ continua y compacta. Si
$x\ne\mathcal Kx$ para todo $x\in\partial\Omega$, está definido el grado de
Leray--Schauder de $I-\mathcal K$. Conserva las propiedades esenciales del
grado finito-dimensional. Si
\[
 \deg_{LS}(I-\mathcal K,\Omega,0)\neq0,
\]
entonces $\mathcal K$ tiene un punto fijo en $\Omega$.

Para el operador \eqref{eq:guide-caputo-fixed-operator}, tómese
$X=C([t_0,T],\R^n)$. Si $f$ es continua y está acotada por $M$ sobre una región
que contiene las funciones consideradas, entonces
\begin{equation}
 \|\mathcal Kx-x_0\|_\infty
 \leq \frac{M(T-t_0)^q}{\Gamma(q+1)}.
 \label{eq:guide-caputo-operator-bound}
\end{equation}
Para verificar equicontinuidad, sean \(t_0\le s\le t\le T\).
Al separar la integral en \([t_0,s]\) y \([s,t]\), el núcleo decreciente
\(u^{q-1}\) produce
\begin{align*}
 \|\mathcal Kx(t)-\mathcal Kx(s)\|
 &\le\frac M{\Gamma(q)}\left[
 \int_{t_0}^{s}\bigl((s-\tau)^{q-1}-(t-\tau)^{q-1}\bigr)\,d\tau
 +\int_s^t(t-\tau)^{q-1}\,d\tau\right]\\
 &=\frac M{\Gamma(q+1)}
 \bigl[(s-t_0)^q-(t-t_0)^q+2(t-s)^q\bigr]\\
 &\le\frac{2M}{\Gamma(q+1)}|t-s|^q.
\end{align*}
Esta cota es uniforme respecto de \(x\). Junto con
\eqref{eq:guide-caputo-operator-bound}, permite aplicar Arzelà--Ascoli
y obtener compacidad relativa de la imagen. La continuidad de
\(\mathcal K\) también debe comprobarse: si \(x_j\to x\) uniformemente
y sus valores pertenecen a un compacto de estados, la continuidad uniforme
de \(f\) en ese compacto implica \(f\circ x_j\to f\circ x\)
uniformemente. Integrar el núcleo da entonces
\[
 \|\mathcal Kx_j-\mathcal Kx\|_\infty
 \le\frac{(T-t_0)^q}{\Gamma(q+1)}
 \|f\circ x_j-f\circ x\|_\infty\longrightarrow0.
\]

Un cálculo concreto del grado proporciona existencia local sin exigir
Lipschitz. Sea \(f\) continua en \(\overline B(x_0,R)\), con
\(\|f\|\le M\), y tome \(T>t_0\) tan próximo a \(t_0\) que
\(M(T-t_0)^q/\Gamma(q+1)<R\). En
\(\Omega=\{x\in X:\|x-x_0\|_\infty<R\}\), considere la homotopía
compacta
\[
 \mathcal H_\lambda x=x_0+\lambda(\mathcal Kx-x_0),
 \qquad 0\le\lambda\le1.
\]
Si \(x=\mathcal H_\lambda x\), la cota anterior fuerza
\(\|x-x_0\|_\infty<R\). Ningún punto fijo alcanza
\(\partial\Omega\), por lo que la normalización y la invariancia
homotópica dan
\[
 \deg_{LS}(I-\mathcal K,\Omega,0)
 =\deg_{LS}(I-x_0,\Omega,0)=1.
\]
Existe, por tanto, una solución del PVI en \([t_0,T]\). La continuidad
por sí sola no garantiza su unicidad.

Este razonamiento es compatible con la memoria porque trabaja directamente
con funciones completas. Sin embargo:

\begin{itemize}
 \item el punto fijo de $\mathcal K$ es una trayectoria en $[t_0,T]$, no un
 punto fijo del campo $f$;
 \item existencia en un horizonte finito no implica atracción asintótica;
 \item el grado no prueba caos, ocultedad ni coexistencia de cuencas;
 \item una continuación en $q$ requiere cotas a priori uniformes y una
 homotopía que no reinicie inadvertidamente la memoria.
\end{itemize}

\subsection{Relación con los métodos de localización y continuación}

\begin{table}[H]
\centering
\caption{Alcance de las formulaciones topológicas.}
\label{tab:topological-method-boundary}
\small
\begin{tabularx}{\textwidth}{@{}L{3.0cm}YY@{}}
\toprule
Objeto & Conclusión posible & Conclusión que requiere otra prueba \\
\midrule
$f(e)=0$ y grado de $f$ & existencia robusta de algún equilibrio & estabilidad, cuenca u ocultedad \\
$P(z)-z=0$ & existencia de un ciclo & que sea atractivo, oculto o caótico \\
Cierre DF/Machado & semilla armónica aproximada & punto fijo del retorno verdadero \\
Grado de $P-I$ & persistencia de algún retorno fijo & forma global y tamaño de su cuenca \\
$x=\mathcal Kx$ en Caputo & existencia de un PVI en horizonte finito & atractor funcional de largo tiempo \\
Sondeos de cuenca & acceso o no acceso en la muestra declarada & topología completa de la frontera \\
\bottomrule
\end{tabularx}
\end{table}

\ifHAFOPedagogy
\subsection{Problemas de aprendizaje}

\begin{enumerate}[label=\textbf{T\arabic*.}]
 \item Construya el residual de disparo para $u''+u=0$ con
 $u(0)=0$, $u(\pi/2)=1$. Demuestre continuidad y encuentre el cero.
 \item Para un retorno escalar continuo $P:[a,b]\to[a,b]$, pruebe mediante el
 teorema del valor intermedio que $P$ tiene un punto fijo. Relacione el
 argumento con Brouwer en dimensión uno.
 \item Calcule el grado de $F(x)=x^3-x$ en regiones que aíslen cada cero y
 compruebe la propiedad de aditividad.
 \item Formule el residual de retorno del PLL respetando la identificación
 angular. Explique por qué $P-I$ debe escribirse en una carta local o mediante
 una diferencia angular envuelta.
 \item Para el operador \cref{eq:guide-caputo-fixed-operator}, derive
 \cref{eq:guide-caputo-operator-bound}. Indique exactamente dónde se usa que
 $0<q<1$ y que $f$ está acotada.
 \item Diseñe una homotopía entre el sistema entero y uno fraccionario. Liste
 las cotas uniformes que faltaría probar antes de invocar grado o continuación.
\end{enumerate}

\paragraph{Pistas.}
En \textbf{T1}, $u(t;s)=s\sin t$ y el residual es afín. Para \textbf{T2},
aplique el teorema del valor intermedio a $P(x)-x$ y use que $P([a,b])$ queda
dentro del mismo intervalo. En \textbf{T3}, aísle primero los ceros y sume los
signos de la derivada. En \textbf{T4}, una resta ordinaria de ángulos no está
bien definida en $S^1$; trabaje en una carta o con una diferencia envuelta. En
\textbf{T5}, la cota sale de integrar exactamente
$(t-\tau)^{q-1}$. En \textbf{T6}, una respuesta completa debe fijar el espacio
de funciones, el terminal inferior, el horizonte, una región acotada y una
cota que sea uniforme en el parámetro de homotopía.
\fi
